\documentclass[11pt]{article}
\usepackage{amsmath,amssymb,booktabs,geometry,graphicx,hyperref,url}
\usepackage{subcaption,placeins}
\hypersetup{hidelinks}
\geometry{margin=1in}
\setlength{\parskip}{0.45em}
\setlength{\parindent}{0pt}
\title{Stage 2: CN/PSOR Solver Validation Audit}
\author{Free-Boundary-Aware Risk Surfaces for American Options}
\date{\today}
\begin{document}
\maketitle

\section{Stage 2: CN/PSOR Solver Validation Audit}

\subsection{Purpose}
This stage audits the finite-difference benchmark before it is allowed to
generate risk surfaces or supervised labels. Crank--Nicolson provides a
time-centered discretization of the Black--Scholes PDE, while projected
successive over-relaxation (PSOR) enforces the American payoff obstacle at each
time step. The audit checks the operator sign, matrix construction, boundary
terms, projection rule, analytic validation cases, shape diagnostics, grid and
domain sensitivity, and runtime.

\subsection{Method}
The solver uses a uniform spot grid on \([0,S_{\max}]\) and a
time-to-maturity grid \(\tau\in[0,T]\). The default smoke grid is \(M=N=120\),
with refinement at \(80,120,160,\) and \(240\). The validation domain uses
\(S_{\max}=4K\), with sensitivity checks at \(3K,4K,\) and \(5K\). The
grid-240 solution is an internal engineering reference rather than an analytic
truth surface.

Validation cases are:
\begin{enumerate}
  \item European put with projection disabled, compared against the Black--Scholes closed form.
  \item No-dividend American call, compared against the European call theorem.
  \item American put LCP diagnostics.
  \item Dividend-paying American call LCP diagnostics and boundary-condition audit.
\end{enumerate}

\subsection{Mathematical Formulation}
The discrete spatial operator approximates
\[
\mathcal{L}U=\frac{1}{2}\sigma^2S^2U_{SS}+(r-q)SU_S-rU.
\]
At interior node $S_i$, the finite-difference coefficients are
\[
a_i=\frac{1}{2}\sigma^2S_i^2/\Delta S^2-\frac{(r-q)S_i}{2\Delta S},
\quad
b_i=-\sigma^2S_i^2/\Delta S^2-r,
\quad
c_i=\frac{1}{2}\sigma^2S_i^2/\Delta S^2+\frac{(r-q)S_i}{2\Delta S}.
\]
Thus $L_h$ has lower, diagonal, and upper entries $(a_i,b_i,c_i)$, and the Crank--Nicolson matrices are
\[
A=I-\frac{\Delta\tau}{2}L_h,\qquad B=I+\frac{\Delta\tau}{2}L_h.
\]
The American LCP at each time step is
\[
AU^{n+1}\ge B U^n,\qquad U^{n+1}\ge \Phi,\qquad
(U^{n+1}-\Phi)^\top(AU^{n+1}-BU^n)=0.
\]
PSOR applies the projected update $x_i\leftarrow\max(\Phi_i,x_i+\omega(y_i-x_i))$.

\subsection{From PDE to Linear Complementarity System}
The matrix calculation can be read in four layers. First, the spatial stencil creates the lower, diagonal, and upper coefficients $(a_i,b_i,c_i)$ for the discrete operator $L_h$. Second, Crank--Nicolson splits this operator across two time levels through $A=I-\Delta\tau L_h/2$ and $B=I+\Delta\tau L_h/2$. Third, boundary values at $S=0$ and $S=S_{\max}$ adjust only the first and last interior right-hand-side equations. Fourth, PSOR solves the resulting inequality system by relaxing each node and projecting it back above the payoff obstacle. The new visual gallery is intended to make these array and matrix operations transparent for reproducible implementation review.

\subsection{Implementation Details}
The implementation audited here is in \path{src/solvers/cn_psor.py}. The helper
\texttt{fd\_operator\_matrix} exposes $L_h$ for direct sign-convention testing.
Boundary conditions are defined in \path{src/solvers/boundary_conditions.py}.

American put boundary conditions:
\[
U(0,\tau)=K,\qquad U(S_{\max},\tau)=0.
\]
Dividend-paying American call boundary conditions:
\[
U(0,\tau)=0,\qquad
U(S_{\max},\tau)=\max(S_{\max}-K,\; S_{\max}e^{-q\tau}-Ke^{-r\tau},\;0).
\]
For European validation with projection disabled, the put lower boundary uses $Ke^{-r\tau}$ instead of $K$.

Rannacher smoothing is not active in the current solver. The configuration field \texttt{rannacher\_steps} is present but not implemented, so near-maturity Greek smoothing should not be claimed from this audit.

\subsection{Validation and Diagnostics}
Unit tests added or run for this audit verify:
\begin{itemize}
  \item the finite-difference operator sign for a linear test function away from the truncated upper boundary;
  \item the Crank--Nicolson matrix signs $A=I-\Delta\tau L_h/2$ and $B=I+\Delta\tau L_h/2$;
  \item American put boundary conditions;
  \item dividend-paying American call boundary conditions;
  \item heterogeneous CSV output for audit tables.
\end{itemize}

The validation experiment writes:
\begin{itemize}
  \item \texttt{validation\_metrics.csv};
  \item \texttt{diagnostics\_table.csv};
  \item \texttt{grid\_refinement.csv};
  \item \texttt{domain\_sensitivity.csv};
  \item \texttt{boundary\_conditions.csv};
  \item \texttt{validation\_summary.json}.
\end{itemize}

\subsection{Results}
On the smoke grid \(M=N=120\), the European put closed-form maximum error is
\(2.782\times10^{-4}\), and the no-dividend American call maximum error is
also \(2.782\times10^{-4}\). These values reproduce the earlier smoke test,
but they are not presented as publication-grade accuracy.

\begin{table}[h]
\centering
\begin{tabular}{lrrrr}
\toprule
Case & Grid & Max error & RMSE & Runtime (s)\\
\midrule
European put & 80 & $6.2445\times10^{-4}$ & $1.6077\times10^{-4}$ & 0.502\\
European put & 120 & $2.7817\times10^{-4}$ & $7.1401\times10^{-5}$ & 0.703\\
European put & 160 & $1.5647\times10^{-4}$ & $4.0166\times10^{-5}$ & 1.176\\
European put & 240 & $6.9375\times10^{-5}$ & $1.7858\times10^{-5}$ & 2.845\\
No-dividend American call & 80 & $6.2445\times10^{-4}$ & $1.6077\times10^{-4}$ & 0.209\\
No-dividend American call & 120 & $2.7817\times10^{-4}$ & $7.1401\times10^{-5}$ & 0.488\\
No-dividend American call & 160 & $1.5647\times10^{-4}$ & $4.0166\times10^{-5}$ & 0.875\\
No-dividend American call & 240 & $6.9375\times10^{-5}$ & $1.7858\times10^{-5}$ & 2.219\\
\bottomrule
\end{tabular}
\caption{Closed-form and theorem-based validation under grid refinement. Errors decrease with refinement, but the grid-240 values remain engineering-scale evidence rather than a universal accuracy certificate.}
\end{table}

\begin{table}[h]
\centering
\resizebox{\textwidth}{!}{%
\begin{tabular}{lrrrrr}
\toprule
Case & Obstacle & Complementarity & Monotonicity & Convexity & Max PSOR\\
\midrule
American put & 0 & $4.0088\times10^{-12}$ & 0 & $2.22\times10^{-16}$ & 12\\
No-dividend American call & 0 & $3.1344\times10^{-11}$ & $2.79\times10^{-17}$ & $4.03\times10^{-10}$ & 18\\
Dividend American call & 0 & $5.2696\times10^{-12}$ & $9.44\times10^{-14}$ & $4.50\times10^{-13}$ & 13\\
\bottomrule
\end{tabular}
}
\caption{Obstacle, complementarity, shape, and PSOR diagnostics on the smoke grid. These discrete checks assess feasibility and residual consistency; they do not replace the separate grid and domain audits.}
\end{table}

Domain sensitivity is stable at this audit level. The American put initial boundary is $0.8K$ for $S_{\max}=3K,4K$ and $0.7917K$ for $S_{\max}=5K$. The dividend-paying call initial boundary is $1.35K$ for $S_{\max}=3K$ and $1.3333K$ for $S_{\max}=4K,5K$. The movement is at the scale of the grid spacing and should be refined before final paper tables.

% Auto-generated by experiments/19_report_visualization_enrichment.py
\section*{Visual Evidence Gallery}
\begin{figure}[!htbp]
\centering
\begin{subfigure}[t]{0.48\linewidth}
\centering
\includegraphics[width=\linewidth]{results/figures/report_enrichment/stage02_solver_validation_tridiagonal_operator.pdf}
\caption{The discretized Black Scholes operator has the expected sparse tridiagonal structure.}
\end{subfigure}
\hfill
\begin{subfigure}[t]{0.48\linewidth}
\centering
\includegraphics[width=\linewidth]{results/figures/report_enrichment/stage02_solver_validation_fd_stencil_validation.pdf}
\caption{The validation audit checks that the implementation uses the intended three point stencil.}
\end{subfigure}
\end{figure}

\begin{figure}[!htbp]
\centering
\begin{subfigure}[t]{0.48\linewidth}
\centering
\includegraphics[width=\linewidth]{results/figures/report_enrichment/stage02_solver_validation_coefficient_sign_audit.pdf}
\caption{The coefficient diagram helps connect operator sign conventions to implemented matrix rows.}
\end{subfigure}
\hfill
\begin{subfigure}[t]{0.48\linewidth}
\centering
\includegraphics[width=\linewidth]{results/figures/report_enrichment/stage02_solver_validation_annotated_lh_matrix.pdf}
\caption{The labeled operator matrix makes the lower diagonal and upper entries explicit for implementation review.}
\end{subfigure}
\end{figure}

\begin{figure}[!htbp]
\centering
\begin{subfigure}[t]{0.48\linewidth}
\centering
\includegraphics[width=\linewidth]{results/figures/report_enrichment/stage02_solver_validation_cn_a_b_matrices.pdf}
\caption{The Crank--Nicolson matrix pair makes the A and B sign convention explicit for the Stage 2 implementation audit.}
\end{subfigure}
\hfill
\begin{subfigure}[t]{0.48\linewidth}
\centering
\includegraphics[width=\linewidth]{results/figures/report_enrichment/stage02_solver_validation_rhs_boundary_terms.pdf}
\caption{The right hand side diagram shows why boundary conditions must be audited together with matrix construction.}
\end{subfigure}
\end{figure}

\begin{figure}[!htbp]
\centering
\begin{subfigure}[t]{0.48\linewidth}
\centering
\includegraphics[width=\linewidth]{results/figures/report_enrichment/stage02_solver_validation_psor_projection_step.pdf}
\caption{The projected update enforces the payoff obstacle after each relaxed candidate update.}
\end{subfigure}
\hfill
\begin{subfigure}[t]{0.48\linewidth}
\centering
\includegraphics[width=\linewidth]{results/figures/report_enrichment/stage02_solver_validation_lcp_matrix_system.pdf}
\caption{The LCP diagram connects the CN matrices to obstacle and complementarity diagnostics.}
\end{subfigure}
\end{figure}

\begin{figure}[!htbp]
\centering
\begin{subfigure}[t]{0.48\linewidth}
\centering
\includegraphics[width=\linewidth]{results/figures/report_enrichment/stage02_solver_validation_time_step_flow.pdf}
\caption{Time marching starts from the nonsmooth payoff and builds the American value surface step by step.}
\end{subfigure}
\hfill
\begin{subfigure}[t]{0.48\linewidth}
\centering
\includegraphics[width=\linewidth]{results/figures/report_enrichment/stage02_solver_validation_validation_ladder.pdf}
\caption{Validation proceeds from analytic checks to inequality and shape diagnostics before downstream experiments.}
\end{subfigure}
\end{figure}

\begin{figure}[!htbp]
\centering
\begin{subfigure}[t]{0.48\linewidth}
\centering
\includegraphics[width=\linewidth]{results/figures/report_enrichment/stage02_solver_validation_smoke_error_comparison.pdf}
\caption{Closed-form European put and theorem-based no-dividend call checks quantify smoke-grid error before downstream classical experiments.}
\end{subfigure}
\hfill
\begin{subfigure}[t]{0.48\linewidth}
\centering
\includegraphics[width=\linewidth]{results/figures/report_enrichment/stage02_solver_validation_obstacle_complementarity.pdf}
\caption{Complementarity residuals remain near solver tolerance across the smoke cases; obstacle violation is zero for every case and is stated explicitly rather than shown as an empty bar chart.}
\end{subfigure}
\end{figure}

\FloatBarrier


\subsection{Failure Modes and Limitations}
The remaining error is most plausibly due to uniform-grid resolution near the payoff kink, finite-domain truncation, and Crank--Nicolson time stepping from nonsmooth terminal data. The refinement table shows decreasing errors from grid 80 to 240, which argues against an operator-sign or boundary-condition implementation error. However, this is still not publication-grade evidence. A final benchmark should use a higher-resolution reference, finer boundary interpolation, and active Rannacher smoothing if near-maturity Greeks are emphasized.

\subsection{Conclusion and Next Step}
The solver-validation gate passes for proceeding to classical boundary, Greek,
risk-surface, grid and domain, and runtime experiments. It does not yet justify
publication-grade accuracy claims. Future reports should continue to attach
obstacle, complementarity, monotonicity, convexity, grid, domain, and runtime
diagnostics to every solver-generated surface.

\end{document}
